\section{Exact cover existence before ambient indexing}
\label{sec:cover-oracle}

The integrated shared Pachner engine removes restarted trace prefixes but
still constructs a complete ambient cover index before visiting the search
root. On the actual-diagram audit of Section~\ref{sec:pachner-native}, this
preparation consumes the local allowance before useful move exploration.
This continuation represents the same existential cover language by exact
on-demand queries on consumed footprints. It preserves connected and
multicomponent covers, the complete bounded family and independent replay.

\subsection{A footprint needs a witness, not the whole index}

For original dual graph $G$, size limit $R$ and component limit $c$, a
footprint $F$ is admitted exactly when
\[
  \exists S:\quad F\subseteq S,\quad |S|\le R,\quad
                          \operatorname{components}(G[S])\le c.
\]
An empty footprint has the empty witness. If $F$ itself satisfies the size
and component limits, it is already a witness. Otherwise the oracle searches
for additional original cells connecting its components. In particular it
must not reject disconnected consumption merely because $c=1$.
A larger connected cover can contain a disconnected footprint.

\begin{proposition}[Rooted completion suffices]
If a permitted cover of $F$ exists, one exists with every induced component
intersecting $F$. The latter covers can be enumerated without duplication
by a rooted reverse search starting from $F$.
\end{proposition}
\begin{proof}
Delete every component of a cover that does not intersect $F$. This retains
all required cells, reduces its size and cannot increase its component
count. For any rooted set $S\supsetneq F$, choose a spanning forest with
one assigned root in $F$ for every tree, for example a deterministic
multi-source breadth-first forest. It has a non-root leaf. Deleting that
leaf leaves every remaining vertex connected to a root, so at least one
non-root deletion preserves the rooted property.
Define the canonical parent by deleting the largest non-root vertex with
that property. Cardinality decreases until $F$ is reached. The deleted
vertex always has a neighbour in its parent, since a component consisting
only of that non-root vertex would violate rootedness. Thus adding frontier
vertices and accepting exactly children whose canonical parent is the
current set generates every rooted extension once.
\end{proof}

The original degree bound is four. For fixed roots $F$, encode any rooted
extension by processing a deterministic breadth-first queue initially
containing those roots. Each processed vertex gives at most four membership
bits for its neighbour ports. Every component meets a root, so the queue
reaches the entire extension. At most $R$ vertices are processed, giving
at most $\sum_{i=0}^{R}2^{4i}=2^{O(R)}$ extensions, independent of an
ambient factor raised to $|F|$. Canonical-parent and component tests cost
polynomial work in $R$ and the source encoding. A fixed query therefore
has $2^{O(R)}\operatorname{poly}(t+B)$ bit cost. No greedy Steiner
approximation or numerical relaxation supplies the answer.

\subsection{The selected witness does not define future eligibility}

A search frame retains its actual consumed original set $F$. Adding an
event's newly consumed originals $A$ asks about $F\cup A$, using the
\emph{original} dual graph. Born cells add no new original identity because
their ancestors entered at creation. An admitted frame stores one covering
witness solely for the source certificate. The next query uses $F$, not
that witness: a later event may require a different cover containing the
larger footprint.

For example on the path $0$--$1$--$2$--$3$ with $R=3,c=1$, consumption
$\{0\}$ has witness $\{0\}$. Later consumption $\{0,2\}$ is still allowed,
with witness $\{0,1,2\}$. Restricting future moves to the first witness
would lose this trace. The new engine does not do so.

The existential predicate is hereditary. If some cover contains a full
trace's footprint, it contains every prefix. It therefore admits exactly
the union language of the indexed engine. Admitted independent-event swaps
have the same final footprint and source marking, and both intermediate
footprints lie in that cover. The same sleep-set commutation argument
remains valid. Endpoint deduplication still does not delete continuation
states with different remaining budgets or obligations.

Witness covers may differ between backends; actual geometry, transported
cochain, incarnation histories and parent trace evidence do not depend on
which covering witness is chosen. Predicate completeness still requires
invariance under marked endpoint/commutation equivalence and independence
from an arbitrarily selected witness cover. A history-sensitive callback
cannot gain a completeness guarantee from changing representation.

\subsection{Completed caches, explicit limits and retained interfaces}

The oracle caches only a completed existence result for a footprint:
one actual witness, or a proved absence. Its global completion-state limit
counts begun rooted-extension candidates. Exhausting that limit or a
cooperative callback interrupts the query; it never returns or caches a
false exclusion. Cheap certified cases, such as a connected footprint
already within the limit, need no extension search.

\code{cover\_backend} accepts \code{auto}, \code{indexed} and \code{oracle}.
The default uses the on-demand backend when no complete-index cap is
requested. An explicit \code{max\_regions} retains its original index-cap
meaning and selects the indexed engine in automatic mode. The two explicit
backends remain differential oracles. \code{max\_oracle\_states} bounds total
completion work across a call, not the ambient cover-family count. These
limits have distinct meanings; an incompatible explicit oracle/index-cap
combination is rejected. The original certificate format and source checker
remain unchanged.

Index counters remain zero on an unmaterialized oracle call; they do not
assert that no logical covers exist. Separate query, cache-hit and extension
state counts describe the actual oracle work. The index implementation and
its cap/measurement regression tests remain available explicitly.

For connected bounded first descent, $R\le3U+3$. Multiplying the existing
$t2^{O(U)}$ geometric trace count by the $2^{O(R)}$ per-footprint completion
cost retains $2^{O(U)}\operatorname{poly}(t+B)$ overall geometry. This replaces
eager preparation with visited-footprint work, rather than improving that
parameterized complexity class. It can expose useful geometric transitions
within an allowance that previously ended during index construction.
For general multicomponent families, no new ambient-independent trace count
is asserted; only each supplied-footprint completion has the stated bound.

\subsection{Actual diagram progress and measured source costs}
All 1,506 maintained tests pass in 275.089 seconds. The 27 focused tests
pass in 2.332 seconds. Their literal small-graph comparison covers every
footprint, size allowance and zero/one/two-component allowance on the selected
graphs. Returned witnesses satisfy the exact requested containment and
component conditions. Tests include disconnected completion, changing the
witness between prefixes, completed negative caches, interrupted/capped
queries, exact indexed work/cap boundaries and independent geometric replay.
A five-node actual trefoil query performs three completed moves with no
ambient index and correctly remains inconclusive.

The 21.817-second audit compares both backends on the preceding genuine
seven- and thirteen-tetrahedron solid tori. Each complete exhaustion preserves
its full marked endpoint proof sequence: twenty-eight and fifty-five entries
agree byte-for-byte after removing only the arbitrary covering-witness field.
The search node counts also agree. All 83 current endpoint proofs replay
independently. The zero-upward bounded exclusions and one-up/two-down strict
descents remain, with fresh Regina confirmations of both original and final
solid tori. Both full source exhaustions require no rooted extension states:
their visited footprints have immediate permitted witnesses.

Thirteen actual diagrams retain all complete enabled recognition verdicts,
and disabled statuses, methods and evidence match the frozen previous
recognizer exactly. Native disc successes are unchanged: the empty-circle
case is positive and receives fresh Regina confirmation, while all twelve
nonempty native calls hit their local allowance. Completed geometric/cochain
transitions nevertheless rise from nine to 141 in the same cohort. Eleven
nonempty diagrams that formerly performed no moves now do so; the remaining
nonempty case rises from nine to thirty-four. This is source exploration
progress, not new knot-recognition coverage.

A separate frozen-source index probe establishes the preparation bottleneck
on three concrete diagrams. With the same 200,000-unit allowance, indexing
fails to finish after 6,092 regions on optimized-positive, 7,079 on genus-one
miss and 7,112 on trefoil. The new five-node pilot completes three moves on
each without an index. Its complete work totals are 123,529, 46,923 and
42,731. The diagram proof query still returns inconclusive; no positive is
inferred from the smaller triangulation or the fact that moves were reached.

The first-descent timings use five randomized paired rounds, forty measured
calls and eight warmups, with identical-baseline A/A controls. Both arms
charge source preparation, completion/index work, their configured complete
positive search, independent strict-descent replay and full serialization.
The preceding indexed source is frozen at \code{ce7fb0737}.

\begin{center}
\small
\begin{tabular}{lrrrrr}
\toprule
Input & previous ms & current ms & old/new & old A/A & new A/A\\
\midrule
descent 1 & 33.207 & 31.849 & 1.033 & 1.010 & 0.999\\
descent 2 & 62.264 & 55.380 & 1.125 & 0.991 & 0.988\\
\bottomrule
\end{tabular}
\end{center}


The source ratios are 1.033 and 1.125. Both searches begin six frames. Work
falls from 6,563 to 5,496 for the smaller source and from 16,854 to 9,665
for the larger. The larger source's A/A controls are about 0.991 and 0.988;
the modest small-source gain is not generalized to a universal speedup.
This improves visited-footprint preparation, without changing the bounded
first-descent complexity class.

The diagram timings complete 160 measured calls and thirty-two warmups.
Both native and complete-recognition rows compare the preceding indexed
stage with the new on-demand stage under the same local cap; the stage is
enabled in both recognition arms. Source/cochain construction, every actual
transition and component query, configured caps, positive replay and full
serialized output are charged. They are distinct from comparing the optional
stage against omitting it altogether.

\textbf{Budgeted native diagram calls, indexed versus on-demand.}
\begin{center}
\small
\begin{tabular}{lrrrrr}
\toprule
Input & previous ms & current ms & old/new & old A/A & new A/A\\
\midrule
empty circle & 9.535 & 9.476 & 1.003 & 1.005 & 0.998\\
optimized positive & 189.580 & 637.641 & 0.294 & 1.015 & 1.000\\
genus one miss & 183.740 & 834.054 & 0.221 & 0.989 & 1.003\\
trefoil & 183.011 & 827.692 & 0.221 & 0.999 & 1.001\\
figure eight & 183.139 & 788.638 & 0.232 & 1.003 & 1.005\\
\bottomrule
\end{tabular}
\end{center}

\textbf{Complete recognition with the optional stage enabled in both arms.}
\begin{center}
\small
\begin{tabular}{lrrrrr}
\toprule
Input & previous ms & current ms & old/new & old A/A & new A/A\\
\midrule
genus one miss & 208.951 & 869.423 & 0.241 & 1.009 & 1.001\\
trefoil & 203.683 & 864.819 & 0.236 & 0.999 & 0.999\\
figure eight & 208.323 & 815.071 & 0.255 & 1.011 & 1.009\\
\bottomrule
\end{tabular}
\end{center}


At the same checkpoint allowance, reaching actual moves costs more wall time
than processing the old eager index. The nonempty native paired ratios are
about 0.29, 0.22, 0.22 and 0.23. Complete enabled recognition ratios are
about 0.24, 0.24 and 0.25: old calls take about 0.20--0.21 seconds, while
current calls take about 0.82--0.87 seconds. The A/A controls are near parity,
so the slower budgeted latency is a real consequence of doing more expensive
geometric work, not claimed as a speed gain. The unchanged source-positive
control is near parity. No additional disc success is demonstrated, and
the optional diagram stage remains off by default.

The useful result is an exact source family that begins exploration without
ambient preprocessing and retains complete fallback/replay guarantees.
Runtime release \code{1519b06f7}, all tests, exact source and actual-diagram
audits, the supplemental index probe and paired timings are retained under
\code{data/cover-oracle-*}. The driver is
\code{causal\_research.cover\_oracle} with \code{audit},
\code{source-benchmark} and \code{diagram-benchmark} modes.
Publication review checks 643 runtime pins, all timing outcomes, all 83
source endpoint proofs, both strict descents and the saved diagram witness
with cover/move/transport producers disabled. General small-budget reach,
terminal structure, global hierarchy accounting and quasi-polynomial
recognition remain unproved.

